\documentclass[12pt, letterpaper]{article}
\usepackage{amsthm, amsmath, amssymb, graphicx, hyperref, url}
\usepackage[utf8]{inputenc}

\newtheorem{thm}{Theorem}
\newtheorem{lemma}[thm]{Lemma}

\theoremstyle{definition}
\newtheorem{ex}{Example}

\theoremstyle{definition}
\newtheorem{defn}[thm]{Definition}

\title{The Irrationality of \(\sqrt{2}\)}
\author{Gordon Chen}
\date{Fall 2026}

\begin{document}
\maketitle

\section{Introduction} \label{intro}
In this short paper we will show that \(\sqrt{2}\) is an irrational number.
In Section \ref{sec2}, we will introduce the concepts of rational and irrational numbers.
In Section \ref{sec3}, we will provide a proof of the fact that \(\sqrt{2}\) is irrational.
Finally, in Section \ref{sec4}, we will discuss what we know about the other numbers that are
known to be irrational.


\section{Preliminaries} \label{sec2}
First, we define the rational numbers.

\begin{defn}
    A number is rational if it is a member of the set
    \[ \mathbb{Q} = \left\{\frac{m}{n} : m, n \in \mathbb{Z} \text{ and } n \ne 0 \right\}.\]
\end{defn}

And now we can define the concept of irrationality.
\begin{defn}
    A real number \(\alpha \in \mathbb{R}\) is irrational if \(\alpha \notin \mathbb{Q}\).
\end{defn}

We will also need the notion of a fraction being in lowest terms.
First we define the greatest common divisor.
\begin{defn}
    The greatest common divisor of integers \(a\) and \(b\), not both zero, is an integer
    \(d\) such that \(d\) divides both \(a\) and \(b\), and \(d\) is greater or equal to every
    common divisor of \(a\) and \(b\). We write \(\gcd(a, b) = d\).
\end{defn}

Now we define a fraction being in lowest terms.
\begin{defn}
    A fraction \(\frac{m}{n}\) where \(m\) and \(n\) are integers and \(n\) is nonzero is in lowest
    terms if the greatest common divisor of \(m\) and \(n\) is 1, that is \(\gcd(m, n) = 1\).
\end{defn}


\section{The Proof} \label{sec3}
\begin{lemma} \label{lemma_div2}
    If \(n^2\) is even for an integer \(n\), then \(n\) is even.

    \begin{proof}
        We proceed by contrapositive. Let \(n\) be an integer and assume that \(n\) is not even,
        so \(n\) is odd.
        Then \(n\) can be written as \(n = 2k + 1\) for some integer \(k\).
        Then \(n^2  = (2k+1)^2 = 4k^2 + 4k + 1 = 2(2k^2 + 2k) + 1\). Since the integers are closed
        under multiplication and addition, then \(2k^2 + 2k\) is an integer. Let \(m = 2k^2 + 2k\).
        Then \(n^2 = 2m + 1\), so \(n^2\) is odd.
        Thus we have shown by that if \(n\) is odd, then \(n^2\) is odd.
        By the contrapositive we then have that if \(n^2\) is even, then \(n\) is even.
    \end{proof}
\end{lemma}

We are finally ready to show our main result.

\begin{thm}
    The real number \(\sqrt{2}\) is irrational.

    \begin{proof}
        Let \(\alpha = \sqrt{2}\). Then \(\alpha^2 = 2\).
        Suppose \(\alpha\) is rational. Then \(\alpha = \frac{m}{n}\), where \(m,n \in \mathbb{Z}\)
        and \(n \ne 0\).
        We can also take \(\frac{m}{n}\) to be in lowest terms, meaning that \(\gcd(m, n) = 1\).
        Then \((\frac{m}{n}) ^ 2 = 2\), so
        \(m^2 = 2n^2\). Then \(m^2\) is even, so by Lemma \ref{lemma_div2}, we have that \(m\) is even.
        Then \(m = 2m'\) for some integer \(m'\), so \(4m'^2 = 2n^2\), so \(2m'^2 = n^2\).
        Then \(n^2\) is even, so again by Lemma \ref{lemma_div2}, we have that \(n\) is even.
        But this is a contradiction because we assumed that \(\frac{m}{n}\) was written in lowest
        terms, that \(m\) and \(n\) share no common divisors greater than 1,
        but since \(m\) and \(n\) are both even, then \(2\) is a common divisor of \(m\) and \(n\).
        Thus by contradiction we have shown that \(\sqrt{2}\) is not rational, so it must be irrational.
    \end{proof}
\end{thm}


\section{Irrational Numbers} \label{sec4}

By generalizing Lemma \ref{lemma_div2} using Euclid's Lemma \cite{euclid_wiki}, we can show that
the square root of any prime number is irrational.
We also know that transcendental numbers \cite{transcendental_wiki} like \(\pi\) and \(e\)
are irrational because transcendental numbers are not algebraic, and rational numbers are algebraic.

The Pythagoreans, working around the 5th and 6th centuries BC \cite{pythag_wiki},
are credited with the discovery that \(\sqrt{2}\) is irrational \cite{sqrt2_wiki}.
Hippasus of Metapontum was said to have drowned at sea for telling those not in the
Pythagorean Society about the irrationality of \(\sqrt{2}\) \cite{hippasus_wiki}.

Later on, in 1844, Joseph Liouville proved the existence of transcendental numbers. Charles Hermite
proved that \(e\) is irrational in 1873, and Ferdinand von Lindemann proved that \(\pi\) is irrational
in 1882 \cite{transcendental_wiki}.


\bibliographystyle{plain}
\bibliography{refs}

\end{document}
